\chapter{How This Curriculum Works (Read This First)}

\begin{brainbox}[What you should walk away with]
The specific exercises are a vehicle for a broader set of skills. By the end of the curriculum, you should have:
\begin{itemize}[nosep]
  \item \textbf{General programming competency} in Unix, Python, and R: enough to parse files, write reusable scripts, manipulate tabular and genomic data, and make publication-quality figures.
  \item \textbf{A mental model of the software stack} that most bioinformatics work depends on: the OS and shell; conda environments and containers; language runtimes and their package ecosystems; domain-specific libraries and pipeline engines. Knowing which layer you are operating at is what lets you diagnose problems instead of flailing.
  \item \textbf{Competent use of AI tools} as part of everyday work: prompting with enough context and constraints to get runnable code, spotting hallucinated APIs, verifying against primary documentation, and treating AI output as a draft rather than an oracle.
  \item \textbf{The ability to ask the right question} before an analysis: what is the biological question, what would a positive result look like, what would a null result look like, what could fool you.
  \item \textbf{Project-structure instincts}: version-controlled code, reproducible environments, ignore-listed raw data, self-explanatory directory layout, journaled progress. Work that someone else (or future you) can clone and re-run.
  \item \textbf{Layer-by-layer evaluation skills}: at each step of a pipeline, the ability to say what should be true, check that it is, and recognize what has gone wrong when it isn't. This is where verification and real understanding live.
\end{itemize}
\end{brainbox}

\begin{brainbox}[Built for how learning actually works]
Traditional curricula front-load reading and back-load doing. This one inverts that order: every session is hands-on first, theory second.

Each session produces something tangible. You run a command, read the output, then read about why it behaves that way.

Jumping around is fine. Each sprint is self-contained enough that you can follow your curiosity without losing the thread.
\end{brainbox}


\section{The Three Rules}

\textbf{\color{teal}Rule 1: The 25-Minute Sprint.} Every study session is broken into 25-minute sprints (Pomodoro). Each sprint has one clear goal and one tangible output. If you finish early, take the break early. If you are in flow, keep going, but never push past 50 minutes without a break.

\textbf{\color{teal}Rule 2: Output Before Input.} Do something before you read about it. Try the command first. Break something. Then read why it broke. You will remember mistakes and fixes much better than anything you passively read.

\textbf{\color{teal}Rule 3: The Checkpoint.} Every sprint ends with a concrete thing you made, produced, or proved works. Screenshot it, commit it, or write one sentence in your journal. This is your proof of progress.


\section{How a Block Works}

The curriculum is divided into 12 \textbf{blocks}. Each block is roughly 2 hours of focused work, organized as a small batch of 25-minute sprints with short breaks between them. Most blocks land around 3--5 sprints; check the chapter for the exact count. How you spread a block across the calendar is up to you: one 2-hour sitting, two 60-minute sittings, or a sprint a day all work. What matters is that you finish the block's sprints before moving on.

\begin{center}
\begin{tabularx}{\textwidth}{l X}
\toprule
\headertext{Within a sitting} & \headertext{What You Do} \\
\midrule
Sprint (25 min) & Hands-on: work through the sprint's goal and produce its checkpoint \\
Break (5 min) & Stand up, walk, look away from the screen \\
Next sprint & Continue, or stop here and pick up another day \\
End of block & Journal one sentence per sprint. Commit to Git. \\
\bottomrule
\end{tabularx}
\end{center}

\textit{Never force yourself past 50 minutes of sprinting without a real break. If you want to do more in a day, take 30+ minutes off before starting the next block.}


\section{Your Toolkit (Set Up in Block 0)}

Before Block 1, spend one sitting setting this up. This is Sprint 0. \textbf{\hyperref[app:tools]{Appendix~A} is the complete install reference.} Every tool named in this curriculum (core infrastructure, command-line utilities, bioinformatics packages, R and Python libraries) is listed there with macOS and Linux install commands and official documentation links. Keep it open while you work through the curriculum. For a fast lookup of programming concepts (loops, pipes, recursion, regex, git verbs), see \hyperref[app:concepts]{Appendix~B}; for a curated reading list and a block-by-block overview, see \hyperref[app:reading]{Appendix~C}.

\begin{tipbox}[If you are brand new to the terminal]
On macOS, open \texttt{Terminal.app} (Cmd-Space, type ``terminal''). On Linux, any terminal emulator works (GNOME Terminal, Konsole, xterm). On Windows, install WSL2 first (\url{https://learn.microsoft.com/windows/wsl/install}) and do everything inside the Ubuntu shell it gives you --- the rest of this curriculum assumes a Unix-like shell. The \texttt{\$} symbol at the start of a command is the shell prompt; do not type it.
\end{tipbox}

\begin{center}
\begin{tabularx}{\textwidth}{l l X}
\toprule
\headertext{What} & \headertext{Why} & \headertext{How} \\
\midrule
Terminal & Everything runs here & Linux/macOS terminal; or WSL2 on Windows --- see \hyperref[app:tools]{Appendix~A} \\
Miniforge & Package management without pain & \url{https://github.com/conda-forge/miniforge} --- see \hyperref[app:tools]{Appendix~A} \\
VS Code & Editor with terminal built in & \url{https://code.visualstudio.com} (or your favorite editor) \\
Git + GitHub & Version control \& your portfolio & \url{https://git-scm.com}; make a GitHub account now \\
Claude or ChatGPT & AI tutor always available & \url{https://claude.ai} (free tier works) \\
A journal.md file & One sentence per day = progress proof & Create it in your Git repo right now \\
\bottomrule
\end{tabularx}
\end{center}

\begin{winbox}[Sprint 0 Checkpoint]
Install everything. Open your terminal. Create a local Git repo called \texttt{bioinfo-journey} with a \texttt{journal.md} file inside it. You will push this repo to GitHub in \hyperref[chap:unix]{Block 2, Sprint 8} once you have an SSH key set up; for now, just initialize it locally with \texttt{git init} and make your first commit.
\end{winbox}


\section{Interleaving: Why This Curriculum Mixes Things Up}

\begin{brainbox}[Why not just learn Unix, then Python, then R in order?]
Research consistently shows that retention improves when topics are interleaved (mixed) rather than massed (all-Unix-then-all-Python).

Each block below has a primary focus, but deliberately connects to other skills. By Block 4 you are already using Python, Unix, and statistics together.

If you get stuck on one topic, switch to a different sprint in the same block. Coming back later with fresh eyes is not quitting; it is how interleaving works.
\end{brainbox}


\section{How to Use AI as Your Study Partner}

AI tools are effective study aids. Treat them as a resource you can query any time for explanations, debugging, and code review.

\begin{itemize}[nosep]
  \item \textbf{Stuck on an error?} Paste the full error message into Claude or ChatGPT. Do not describe it; paste it.
  \item \textbf{Confused by a concept?} Ask: ``Explain [concept] assuming I already know [thing I do know] but not [thing I don't].''
  \item \textbf{Want to check understanding?} Ask: ``I think [my understanding]. Is that right? What am I missing?''
  \item \textbf{Need a break from reading?} Ask: ``Give me a 5-minute hands-on exercise to practice [concept].''
  \item \textbf{Feeling overwhelmed?} Ask: ``What is the single most important thing I should understand about [topic]?''
\end{itemize}

\begin{warnbox}[The one rule about AI]
Always run AI-generated code before trusting it. AI tools hallucinate function names, invent packages, and mix up API versions. Test first, trust after.
\end{warnbox}
